\SourceLeaf{638}
\begin{center}ON USING THE COMPARATOR.\end{center}
\begin{center}[1.]\end{center}
Three microscopes set at nearly equal distances apart are placed in succession over intervals of a scale that are nearly equal to those distances.

In these successive trials, the graduation marks on the scale fall beyond the lines of sight of the microscopes by the following amounts
\[
\begin{array}{ccc}
a^0&b^0&c^0\\
a'&b'&c'\\
a''&b''&c''\\
a'''&b'''&c'''\\
&\text{etc.}&
\end{array}
\]
We denote the errors in these quantities by $x^0,y^0,z^0;\ x',y',z';\ x''$, etc. If one sets
\[
\begin{aligned}
a'-a''-b^0+b''+c^0-c'&=e^0\\
a''-a'''-b'+b'''+c'-c''&=e'
\end{aligned}
\]
and so on, the error equations are
\[
\begin{aligned}
x'-x''-y^0+y''+z^0-z'&=e^0\\
x''-x'''-y'+y'''+z'-z''&=e'\ \text{and so on}
\end{aligned}
\]
It follows that the most plausible values of the errors take the form
\[
\begin{array}{lll}
x^0=0&y^0=-h^0&z^0=h^0\\
x'=h^0&y'=-h'&z'=h'-h^0\\
x''=h'-h^0&y''=-h''+h^0&z''=h''-h'\\
x'''=h''-h'&y'''=-h'''+h'&z'''=h'''-h''\\
\text{and so on}&&
\end{array}
\]
\SourceLeaf{639}
and the auxiliary quantities $h^0,h'$, etc., depend on the equations
\[
\begin{aligned}
6h^0-2h'-h''&=e^0\\
-2h^0+6h'-2h''-h'''&=e'\\
-h^0-2h'+6h''-2h'''-h''''&=e''\\
-h'-2h''+6h'''-2h''''-h^{\mathrm V}&=e'''\\
-h''-2h'''+6h''''-2h^{\mathrm V}-h^{\mathrm {VI}}&=e''''\\
\text{etc.}&
\end{aligned}
\]

\par\Needspace{5\baselineskip}
\begin{center}[2.]\end{center}
Arrangement of the length comparisons for determining an interval on a graduated scale in units of the entire scale.

\textit{Example.} The graduation marks on the scale were numbered from 0 to 840. The aim was to determine the interval from 0 to 87 in units of the whole scale from 0 to 840. — As an abbreviation, let
\[
0\ .\ 87\ .\ 174=\alpha
\]
mean that comparison of the length from 0 to 87 with the length from 87 to 174 showed the former to exceed the latter by $\alpha$ divisions of the micrometer, and so on.
\begin{center}Arrangement:\end{center}
\[
\begin{array}{rcrcrcl}
0&.&87&.&174&=&\alpha\\
0&.&174&.&348&=&\beta\\
0&.&348&.&696&=&\gamma\\
840&.&696&.&552&=&\delta\\
840&.&552&.&264&=&\varepsilon\\
0&.&264&.&528&=&\zeta\\
840&.&528&.&216&=&\eta\\
0&.&216&.&432&=&\theta\\
840&.&432&.&24&=&\iota\\
0&.&24&.&48&=&\kappa\\
0&.&48&.&96&=&\lambda\\
0&.&96&.&192&=&\mu\\
0&.&192&.&384&=&\nu\\
0&.&384&.&768&=&\pi\\
840&.&768&.&696&=&\rho
\end{array}
\]

\SourceLeaf{640}
\textit{Calculation.} Denote the full length of the scale by $840'$, the length sought (from 0 to 87) by $87'$, the length from 87 to 174 by $174'-87'$, and so on. This gives the following equations:
\[
\begin{aligned}
2\ .\ 87'-\alpha&=174'\\
2\ .\ 174'-\beta&=348'\\
2\ .\ 348'-\gamma&=696'\\
2\ .\ 696'-840'+\delta&=552'\\
2\ .\ 552'-840'+\varepsilon&=264'\\
2\ .\ 264'-\zeta&=528'\\
2\ .\ 528'-840'+\eta&=216'\\
2\ .\ 216'-\theta&=432'\\
2\ .\ 432'-840'+\iota&=24'\\
2\ .\ 24'-\kappa&=48'\\
2\ .\ 48'-\lambda&=96'\\
2\ .\ 96'-\mu&=192'\\
2\ .\ 192'-\nu&=384'\\
2\ .\ 384'-\pi&=768'\\
2\ .\ 768'-840'+\rho&=696'
\end{aligned}
\]
The result is
\[
\begin{aligned}
696'&=-840'+\rho-64.840'+64\iota-256.840'+256\eta-1024.840'\\
&\quad-2\pi-128\theta-512\zeta+1024\varepsilon\\
&\quad-4\nu-2048.840'\\
&\quad-8\mu+2048\delta\\
&\quad-16\lambda+4096.696'\\
&\quad-32\kappa
\end{aligned}
\]
Thus, if $840'=L$, $696'=y$, $87'=z$,
\SourceLeaf{641}
\[
\begin{aligned}
3393L&=4095y+\rho\\
&\qquad-2\pi\\
&\qquad-4\nu\\
&\qquad-8\mu\\
&\qquad-16\lambda\\
&\qquad-32\kappa\\
&\qquad+64\iota\\
&\qquad-128\theta\\
&\qquad+256\eta\\
&\qquad-512\zeta\\
&\qquad+1024\varepsilon\\
&\qquad+2048\delta
\end{aligned}
\]
\[
z=\tfrac18y+\tfrac12\alpha+\tfrac14\beta+\tfrac18\gamma
\]
Thus the required value of $z=87'$ in units of the full scale $L=840'$ is obtained if the micrometer is checked and the value of its divisions is thereby found to be uniform and is given in units of the distance between two graduation marks of the scale set according to the arrangement above (for example, in units of the length from 87 to 96).

